% ЗӨВХӨН БҮТЭН НООРОГТ (main.tex) хэвлэгдэнэ. Багшид үзүүлэх ноорогт (main_teacher.tex) орохгүй.
% 1.3 Судалгааны асуулт ба хэмжигдэх шалгуур
\section{Судалгааны асуулт ба хэмжигдэх шалгуур}
\label{sec:rq}

Дөрөвдүгээр зорилтыг хэмжүүртэй холбохын тулд хоёр судалгааны асуултыг тодорхойлов.
\begin{enumerate}
    \item Системийн хайлт ресепшнд асуудаг бодит асуултын хэдэд нь зөв эх сурвалжийг олдог вэ?
    \item Хариулт нь сургуулийн мэдээлэлд байхгүй асуултын хэдэд нь систем хандах алба, өрөөг зөв заадаг вэ?
\end{enumerate}

Эхний асуултын хэмжүүр нь зөв эх сурвалж хайлтын эхний нэг болон эхний гурван үр дүнд орсон асуултын тоо юм. Асуултын хүрээгээр нэмэлт оноо өгөх тохиргооны (\ref{sec:rag-method}-р хэсэг) нөлөөг ялгаж харахын тулд энэ тоог нэмэлт оноо өгөөгүй үеийн тоотой нэг хүснэгтэд тайлагнана. Хоёр дахь асуултын хэмжүүр нь хариулт байхгүй асуултаас зөв алба, өрөөг заасан асуултын тоо юм. Мөн алдааны шалтгааныг тайлбарлах туслах хэмжүүр болгож асуултын хүрээг зөв тодорхойлсон тоог тайлагнана. Асуулт бүрийн зөв хариуг системийн хариултыг харахаас өмнө тэмдэглэнэ. Үнэлгээний аргыг \ref{sec:exp-tree}-р хэсэгт тодорхойлов.
% 2026.11.15 (3), хэрэглэгч: суурь тоог (нэмэлт оноогүй) хамт тайлагнана — «хоёр арга харьцуулах
% туршилт» гэж нэрлэхгүй, нэг хүснэгтийн нэг багана.
